\scanpage{46}
\[
\deleted{\varphi:X(K)\times M_K\longrightarrow\mathbb R}
\]
\deleted{est dite \uline{loc.} $M_K$-\uline{bornée}, \inserted{ou \uline{admissible}} si elle est \inserted{$M_K$-}bornée sur les parties $M_K$-bornées, \inserted{i.e. si pour tt tel $E$} Noter que \deleted{\unsure{une partie de}} \inserted{pour \unsure{tte|tt}} \deleted{partie finie $F$} \inserted{$P\in X(K)$,} \unsure{de} \deleted{$X(\overline K)\times X(\overline K)$}, $\{P\}\times M_K$ est $M_K$-borné, donc si $\varphi$ est loc. $M_K$-bornée, alors \deleted{\unsure{la}} \inserted{pour $P\in X(K)$,} $\varphi(P,\cdot):M_K\longrightarrow\mathbb R$ est un \deleted{$M_{\overline K}$} \deleted{$M_{\overline K}$-diviseur} $M_K$-diviseur. Donc une $\varphi$ admissible s'interprète aussi comme une application}
\[
\deleted{\widetilde{\varphi}:X(K)\longrightarrow\operatorname{Div}(M_K)}
\]
\deleted{qui a \inserted{tout au moins} la propriété suivante : Si $F$ est une partie $M_K$-bornée de $X(K)$ [i.e. telle que $F\times M_K$ soit $M_K$-borné], \uline{alors} l'ens des $\widetilde{\varphi}(P)$ $(P\in F)$ est \uline{borné} dans l'ens des diviseurs de $M_K$ [\uline{considéré comme ens ordonné}].}

\margin{\deleted{?}}
\deleted{\uline{Cette condition est-elle suffisante} ??}

\deleted{\deleted{\unsure{ou|on}} Une fonction $\varphi:X(K)\times M_K\longrightarrow\mathbb R$ qui est $M_K$-bornée est aussi dite \uline{\unsure{équivalente à zéro}} (elle est alors admissible)}
% archive sheet number: 45
